\documentclass[11pt,reqno,a4paper]{amsart}

\usepackage[T1]{fontenc}
\usepackage{lmodern}
\usepackage{amsmath,amssymb,mathtools}
\usepackage{booktabs}
\usepackage{array}
\usepackage{enumitem}
\usepackage{microtype}
\usepackage[hidelinks]{hyperref}

\newtheorem{theorem}{Theorem}[section]
\newtheorem{proposition}[theorem]{Proposition}
\newtheorem{lemma}[theorem]{Lemma}
\newtheorem{corollary}[theorem]{Corollary}
\theoremstyle{definition}
\newtheorem{definition}[theorem]{Definition}
\theoremstyle{remark}
\newtheorem{remark}[theorem]{Remark}

\newcommand{\comp}[1]{\overline{#1}}
\newcommand{\K}{K}
\newcommand{\cF}{\mathcal F}
\newcommand{\cP}{\mathcal P}
\newcommand{\emptycert}{+\infty}
\newcolumntype{L}[1]{>{\raggedright\arraybackslash}p{#1}}

\title[The nine-color case]{The Nine-Color Case of an
Erd\H{o}s--Gy\'arf\'as Balanced-Coloring Problem}
\author{Robert Sneiderman}
\date{Preprint, 21 July 2026}
\subjclass[2020]{05C15, 05C35, 05D10, 68V20}
\keywords{edge coloring, extremal graph theory, computer-assisted proof,
LRAT certificate}

\begin{document}
\raggedbottom
\emergencystretch=2em
\hypersetup{
  pdftitle={The Nine-Color Case of an Erdos--Gyarfas Balanced-Coloring Problem},
  pdfauthor={Robert Sneiderman}
}

\begin{abstract}
Erd\H{o}s and Gy\'arf\'as asked whether every $r$-coloring of
$E(K_{r^2+1})$ contains $r+1$ vertices on which some color is absent.
We prove the assertion for $r=9$: every nine-coloring of $E(K_{82})$
has ten vertices whose induced edges omit a color. The proof combines a
colored extremal recursion with two terminal results. The first proves
$P_3(26)\ge121$ by exact rational certificates, structural reductions,
and solver-free finite-state searches. The second proves that the
corresponding family on 27 vertices is empty; a human reduction leaves
50 finite instances, each refuted by a replayed LRAT certificate. A
short full-color argument then proves $P_4(37)\ge192$, which makes every
cell of the outer packing recursion strict. The source release contains
independent catalog reconstruction, semantic verifiers, corruption tests,
manifests, and replay instructions. This fixed case does not settle the
problem for arbitrary $r$.
\end{abstract}

\maketitle

\begin{center}
\fbox{\parbox{0.91\textwidth}{\centering
\textbf{Verification status.}
The solver-free proof, full implication chain, independent
101,880-state CaDiCaL reconstruction, and clean public-package replay pass
locally. External mathematical review remains pending.}}
\end{center}

\section{Introduction}\label{sec:introduction}

Erd\H{o}s and Gy\'arf\'as posed the following question
\cite{ErdosGyarfas1999}. If the edges of $K_{r^2+1}$ are colored with
$r$ colors, must some set of $r+1$ vertices omit a color on its induced
edges? The assertion is immediate for some small parameters, but no
general proof is known. It is listed as Erd\H{o}s Problem 617.

We establish the nine-color case.

\begin{theorem}\label{thm:main}
Every map
\[
  \chi:E(K_{82})\longrightarrow[9]
\]
has a set $S\subseteq V(K_{82})$ of order ten such that
\[
  \chi(E(K_{82}[S]))\ne[9].
\]
\end{theorem}

The argument is by contradiction. For each color $i$, let $G_i$ contain
the edges of color $i$. If every ten-set sees all colors, then
\[
  \alpha(G_i)\le9
  \quad\text{and}\quad
  e(G_i[S])\le37\quad(|S|=10).
\]
More importantly, the eight graphs of the other colors partition
$E(\comp{G_i})$. This simultaneous decomposition gives induced-density
bounds unavailable in a one-color relaxation.

The proof has three parts.
\begin{enumerate}[label=\textup{(\roman*)},leftmargin=2.8em]
\item A colored recursive lower bound turns terminal information into
five forced copies of $K_9$ in a least-color nonneighborhood.
\item Two computer-assisted terminal theorems give
\[
  P_3(26)\ge121,\qquad \cP_3(27)=\varnothing.
\]
\item A human full-color lemma combines these statements to prove
$P_4(37)\ge192$. The resulting outer margins are
$5,4,3,2,3$ in the only formerly deficient row.
\end{enumerate}

The finite work is not treated as an opaque solver call. The order-$26$
proof has exact catalog semantics, rational dual checks, and deterministic
solver-free searches. The order-$27$ proof uses LRAT refutations only
after a human lemma removes 282 of 332 core types. An independent verifier
rebuilds every remaining formula and replays every proof.

\section{Colored induced graphs}\label{sec:colored}

Assume throughout Sections~\ref{sec:colored}--\ref{sec:final} that a
counterexample to Theorem~\ref{thm:main} exists. Fix a target color $i$.
For $n=9q+b$, $0\le b<9$, put
\begin{equation}\label{eq:p9}
 p_9(n)=(9-b)\binom q2+b\binom{q+1}{2}
       =9\binom q2+qb.
\end{equation}
This is the minimum edge count in an $n$-vertex graph with independence
number at most nine.

\begin{lemma}[Full-color induced density]\label{lem:density}
For every $W\subseteq V(K_{82})$,
\[
 e(G_i[W])\le D_9(|W|),\qquad
 D_9(n):=\binom n2-8p_9(n).
\]
\end{lemma}

\begin{proof}
Each nontarget graph $G_j[W]$ has independence number at most nine, so
Tur\'an's theorem gives $e(G_j[W])\ge p_9(|W|)$. The eight nontarget
graphs partition the complement of $G_i[W]$. Summing their lower bounds
proves the claim.
\end{proof}

The values used later are
\begin{equation}\label{eq:D-values}
 D_9(10)=37,\qquad D_9(11)=39,\qquad
 D_9(26)=125,\qquad D_9(27)=135.
\end{equation}

\begin{definition}\label{def:family}
Let $\cF(a,n)$ be the family of graphs $G_i[W]$ inherited from the same
hypothetical coloring, with
\[
 |W|=n,\qquad \alpha(G_i[W])\le a,\qquad \omega(G_i[W])\le8.
\]
If the family is nonempty, let $P_a(n)$ be its minimum edge count.
We also write $\cP_a(n)$ for the family itself and
$\cP_a(n)=\varnothing$ when it is empty.
\end{definition}

The inherited provenance is part of the definition. In particular, every
member and every induced subgraph satisfies Lemma~\ref{lem:density}.

\subsection{The colored recursive bound}

We first justify the terminal thresholds used by the recursion.

\begin{lemma}[Colored layer thresholds]\label{lem:thresholds}
There is no member of $\cF(s,n)$ when
\[
 n\ge9s+T_s,\qquad
 (T_2,T_3,T_4,T_5,T_6)=(0,1,2,4,8).
\]
\end{lemma}

\begin{proof}
For $s=2$, let $H\in\cF(2,n)$ and put $L=\comp H$. The graph $L$ is
triangle-free and $\alpha(L)\le8$, so $\Delta(L)\le8$ and
$e(L)\le4n$. Lemma~\ref{lem:density} gives
$e(L)\ge8p_9(n)$. At $n=18$, equality is forced and $L$ is
8-regular. Since $8>2n/5$, the Andrásfai--Erd\H{o}s--S\'os theorem
makes $L$ bipartite \cite{AndrasfaiErdosSos1974}, contradicting
$\alpha(L)\le8$. For $n>18$, one has $p_9(n)>n/2$, giving an
immediate edge-count contradiction.

Proceed by induction on $s$. If $H\in\cF(s,n)$ and $L=\comp H$, then
every neighborhood in $L$ induces the complement of a member of the
preceding layer. Hence
\begin{equation}\label{eq:threshold-degree}
 \Delta(L)\le9(s-1)+T_{s-1}-1.
\end{equation}
Write $n=9s+t$, $1\le t<9$. For every nontarget color $j$, the graph
$\comp{G_j[W]}$ contains $H$ and is not 9-partite: nine independent
sets of $H$ have total order at most $9s<n$. The
Kang--Pikhurko theorem therefore gives
\[
 e(G_j[W])\ge p_9(n)+s-1.
\]
Summing over the eight nontarget colors and comparing with
\eqref{eq:threshold-degree} gives a contradiction exactly when
\[
 A_st-\widehat B_s>0,
\]
where
\[
 A_s=9(s+1)-2s-T_{s-1}+1,\qquad
 \widehat B_s=9s(s+T_{s-1}-2)-16(s-1).
\]
Thus the first eligible offset is
\[
 T_s=\max\left\{1,
 \left\lfloor\frac{\widehat B_s}{A_s}\right\rfloor+1\right\},
\]
provided it is less than nine. Starting from $T_2=0$ gives
$T_3=1$, $T_4=2$, $T_5=4$, and $T_6=8$. The edge-count margin
increases at every later order: its doubled one-vertex increment is at
least
\[
 16\left\lfloor\frac n9\right\rfloor-
 \bigl(9(s-1)+T_{s-1}-1\bigr)\ge A_s>0.
\]
Thus each contradiction persists beyond its first offset.
\end{proof}

We record the exact recursion used in the final calculation. For an
independence cap $a$, define the Tur\'an floor
\[
 p_a(n)=(a-b)\binom q2+b\binom{q+1}{2},
 \qquad n=aq+b,\quad0\le b<a.
\]
Set
\[
 B(0,0)=0,\qquad B(0,n)=\emptycert\ (n>0),
\]
and
\[
 B(1,n)=
 \begin{cases}
 \binom n2,&n\le8,\\
 \emptycert,&n\ge9.
 \end{cases}
\]
For $a\ge2$, put
\begin{equation}\label{eq:Q}
 Q(a,n)=p_a(n)+
 \begin{cases}
 0,&n\le8a,\\
 \lfloor n/a\rfloor-1,&8a<n\le9a,\\
 \lfloor n/a\rfloor,&n\ge9a+1,
 \end{cases}
\end{equation}
and
\begin{equation}\label{eq:C}
 C(d)=
 \begin{cases}
 \binom{d+1}{2},&d<8,\\
 37,&d=8,\\
 d+\left\lceil\frac{11}{9}\binom d2\right\rceil,&d\ge9.
 \end{cases}
\end{equation}
By Lemma~\ref{lem:thresholds}, the exact colored-layer thresholds used
here are
\[
 T_2=0,\quad T_3=1,\quad T_4=2,\quad T_5=4,\quad T_6=8.
\tag{\ref{eq:C}a}
\]
Whenever $T_a$ is defined and $n\ge9a+T_a$, set
$B(a,n)=\emptycert$. Otherwise set
\begin{equation}\label{eq:R}
 R(a,n)=
 \min_{\substack{0\le d<n\\B(a-1,n-1-d)<\emptycert}}
 \max\left\{
 B(a-1,n-1-d)+C(d),\,
 \left\lceil\frac{nd}{2}\right\rceil
 \right\},
\end{equation}
then put $b(a,n)=\max\{Q(a,n),R(a,n)\}$. If
$b(a,n)>D_9(n)$, set $B(a,n)=\emptycert$; otherwise set
$B(a,n)=b(a,n)$.

\begin{proposition}[Recursive lower bound]\label{prop:recurrence}
Every $F\in\cF(a,n)$ satisfies $e(F)\ge B(a,n)$. A value
$B(a,n)=\emptycert$ certifies that $\cF(a,n)$ is empty.
\end{proposition}

\begin{proof}
Induct on $a$. The cases $a=0,1$ follow from the definitions. The scalar
bound \eqref{eq:Q} is Tur\'an's bound plus the Kang--Pikhurko
nonpartite increment \cite{KangPikhurko2005}. The complement of $F$ is
not $a$-partite when $n>8a$, since one part would be a target clique of
order at least nine. To justify the last extra
unit, consider equality in their construction. If $n\ge9a+2$, its $a$
old parts have total order $n-1>9a$, so one gives a target clique of
order at least ten. If $n=9a+1$, every old part has order nine. For the
special old part $S$, exceptional vertices $x,y$, and nonempty proper
set $A\subset S$, the two ten-sets $S\cup\{x\}$ and $S\cup\{y\}$
have respectively
\[
 \binom92+9-|A|
 \quad\text{and}\quad
 \binom92+|A|
\]
target edges. The cap 37 would require both $|A|\ge8$ and
$|A|\le1$, a contradiction. This proves \eqref{eq:Q}.

Choose a minimum-degree vertex $x$ of $F$, of degree $d$, and let $U$
be its nonneighbor set. Then $F[U]\in\cF(a-1,n-1-d)$. If
$X=e_F(N(x),U)$ and $Y=e_F(N(x))$, minimum degree gives
\[
 X+2Y\ge d(d-1).
\]
For $d<8$ this yields
$d+X+Y\ge\binom{d+1}{2}$. At $d=8$, equality would produce a
target $K_9$, giving one additional edge. For $d\ge9$, averaging the
ten-set cap over the nine-subsets of $N(x)$ gives the last line of
\eqref{eq:C}. The degree-sum bound also gives
$e(F)\ge\lceil nd/2\rceil$. These two estimates yield
\eqref{eq:R}. Lemma~\ref{lem:density} justifies the final emptiness test,
and Lemma~\ref{lem:thresholds} justifies the declared terminal values.
\end{proof}

The terminal results proved below strengthen this recursion by setting
\begin{equation}\label{eq:terminal-inputs}
 B(3,26)\ge121,\qquad B(3,27)=\emptycert.
\end{equation}
A third input, proved in Section~\ref{sec:bridge}, is
\begin{equation}\label{eq:bridge-input}
 B(4,37)\ge192.
\end{equation}

\section{The packing contradiction}\label{sec:packing}

\begin{lemma}[Representative selection]\label{lem:representatives}
Let $Q$ be a target copy of $K_9$ and let $x\notin Q$. At most one
edge from $x$ to $Q$ has the target color. Consequently, from $k$
disjoint target copies of $K_9$ one can choose representatives
target-anticomplete to any fixed target-independent set $I$, provided
$|I|+k\le9$.
\end{lemma}

\begin{proof}
The ten-set $Q\cup\{x\}$ must see all eight nontarget colors. Those
colors can occur only on the nine edges from $x$ to $Q$, so at most one
of those edges has the target color. Choose representatives successively.
Before choosing from the next block, at most eight earlier vertices are
present; each excludes at most one of its nine vertices. At least one
choice remains.
\end{proof}

Let $G$ be a color graph with the fewest edges. Then
\begin{equation}\label{eq:least}
 e(G)\le\frac{1}{9}\binom{82}{2}=369.
\end{equation}
Its minimum degree is at most eight. Indeed, minimum degree at least nine
would force $G$ to be nine-regular. Since $\alpha(G)\le9$, some component
has chromatic number at least ten, contradicting Brooks's theorem unless
it is $K_{10}$ \cite{Brooks1941}. A target $K_{10}$ contradicts the
ten-set cap.

Fix a minimum-degree vertex $v$ of degree $d\le8$, and let $U$ be its
target nonneighbor set. Suppose $j$ disjoint target copies of $K_9$ have
been selected in $U$, and let $R_j$ be the residual. If $R_j$ has no
further $K_9$, representative selection gives
\begin{equation}\label{eq:residual-family}
 R_j\in\cF(8-j,81-d-9j).
\end{equation}
To see the independence bound, an independent $(9-j)$-set in $R_j$
could be extended by one representative from each selected block using
Lemma~\ref{lem:representatives}. The resulting independent nine-set in
$U$, together with $v$, would contradict $\alpha(G)\le9$.

Put $N=N_G(v)$, $X=e_G(N,U)$, and $Y=e(G[N])$. Summing the degree
lower bound over $N$ gives
\[
 X+2Y\ge d(d-1).
\]
Since $Y\le\binom d2$, it follows that $X+Y\ge\binom d2$. The
least-color budget therefore gives
\begin{equation}\label{eq:outer-upper}
 e(R_j)\le
 E_{d,j}:=369-d-\binom d2-36j.
\end{equation}
Therefore
\begin{equation}\label{eq:packing-test}
 B(8-j,81-d-9j)>E_{d,j}
\end{equation}
forces the packing to extend by one block.

\begin{lemma}[Five-block exclusion]\label{lem:five-block}
The graph $G[U]$ cannot contain five disjoint copies of $K_9$.
\end{lemma}

\begin{proof}
Extend five blocks to a maximal packing of $k$ blocks and put
$s=8-k$. The same representative argument gives independence number at
most $s$ in the residual, whose order is
\[
 9s+(9-d).
\]
For $k=5$, this contradicts the colored threshold $T_3=1$. For $k=6$,
the threshold $T_2=0$ applies. If $k=7$, the residual has independence
number at most one and order at least ten, so it is a target clique,
contrary to maximality. If $k\ge8$, eight blocks leave at least one
vertex of $U$ outside them. Lemma~\ref{lem:representatives} selects one
representative from each block target-anticomplete to that vertex. These
nine vertices form an independent set in $G[U]$, and adding $v$ gives
an independent ten-set in $G$. Thus no maximal extension of five blocks
exists.
\end{proof}

\section{The 26-vertex terminal}\label{sec:o26}

\begin{theorem}\label{thm:o26}
Every $H\in\cF(3,26)$ has at least 121 edges.
\end{theorem}

\begin{proof}
Suppose $e(H)\le120$. A vertex of degree at most seven has at least
18 target nonneighbors. Their induced graph belongs to $\cF(2,n)$ for
some $n\ge18$, contrary to Lemma~\ref{lem:thresholds}. Hence
$\delta(H)\ge8$, while averaging gives $\delta(H)\le9$.

The degree-eight branch is excluded by
Lemma~\ref{lem:two-row}, proved below. In the degree-nine branch choose a
degree-nine vertex $v$ and put
\[
 A=N_H(v),\quad B=V(H)\setminus(A\cup\{v\}),\quad
 F=\comp{H[A]},\quad L=\comp{H[B]}.
\]
Thus $|A|=9$ and $|B|=16$. The graph $L$ is triangle-free, so
$e(L)\le64$. Lemma~\ref{lem:density} on $B$ gives
$e(L)\ge8p_9(16)=56$. Hence
\[
 56\le e(L)\le64.
\]
Proposition~\ref{prop:o26-finite} excludes every one of these nine levels.
\end{proof}

\subsection{A common human exclusion}

\begin{lemma}[Seventeen-vertex two-row exclusion]\label{lem:two-row}
Neither of the following configurations exists:
\begin{enumerate}[label=\textup{(\roman*)},leftmargin=2.5em]
\item $|V(H)|=26$, $\delta(H)=8$, and $e(H)\le120$;
\item $|V(H)|=27$, $\delta(H)=9$, and $e(H)\le135$,
\end{enumerate}
where $H\in\cF(3,|V(H)|)$.
\end{lemma}

\begin{proof}
Choose a minimum-degree vertex $v$, put $A=N_H(v)$ and
$B=V(H)\setminus(A\cup\{v\})$, and set
$F=\comp{H[A]}$, $L=\comp{H[B]}$. In both cases $|B|=17$.
The checked core catalog has fourteen types. Every type has a vertex
$z$ such that $L-z$ is bipartite with sides $X,Y$ of order eight and
$d_L(z)=p+q\le5$, where $p$ and $q$ count its neighbors in $X$ and
$Y$. Both are positive. If $h$ is the number of missing $X$--$Y$
edges, then
\begin{equation}\label{eq:h-bound}
 h\le p+q.
\end{equation}

The graph $F$ has an edge $cu$. Put
$S=(N_H(c)\cap B)\cup(N_H(u)\cap B)$. This is a vertex cover of $L$;
otherwise $c,u$ and an uncovered edge of $L$ form an independent
four-set in $H$. Applying $D_9(11)=39$ to
$\{c,u\}\cup X\cup\{z\}$ and to the analogous set using $Y$ gives
\[
 |S\cap X|+\mathbf1_{z\in S}\le p+3,\qquad
 |S\cap Y|+\mathbf1_{z\in S}\le q+3.
\]
If $z\in S$, the independent set $B\setminus S$ forces
\[
 h\ge(6-p)(6-q)>p+q.
\]
If $z\notin S$, the missing neighborhood rectangle and the remaining
independent cross-pairs are disjoint, giving
\[
 h\ge pq+(5-p)(5-q)>p+q.
\]
Both inequalities contradict \eqref{eq:h-bound}. The finite dependency
is only the fourteen-core catalog and its displayed witness.
\end{proof}

\subsection{The nine core levels}

For $a\in A$, write
\[
 D_a=N_H(a)\cap B,\qquad q_a=|D_a|.
\]
The finite reduction uses the following necessary conditions:
\begin{enumerate}[label=\textup{(C\arabic*)},leftmargin=3.2em]
\item $L$ is triangle-free, $\alpha(L)\le8$, and $\Delta(L)\le8$;
\item $F$ is $K_4$-free and $\alpha(F)\le7$;
\item minimum degree gives $q_a\ge d_F(a)$ and column demands
\[
 \sum_{a\in A}\mathbf1_{b\in D_a}\ge d_L(b)-6;
\]
\item if $aa'\in E(F)$, then $D_a\cup D_{a'}$ covers $L$;
\item if $a,a',a''$ form a triangle of $F$, then
$D_a\cup D_{a'}\cup D_{a''}=B$;
\item no row $D_a$ contains an independent eight-set of $L$;
\item the target-edge ledger bounds $\sum_aq_a$ at the relevant level.
\end{enumerate}
Each condition follows from minimum degree, the exclusion of independent
four-sets, or the exclusion of target $K_9$.

\begin{proposition}[Finite order-$26$ exclusion]\label{prop:o26-finite}
No incidence system satisfying \textup{(C1)--(C7)} exists at any core
level $e(L)=56,\ldots,64$.
\end{proposition}

\begin{proof}
Canonical generation lists every eligible core and shell up to
isomorphism. Exact rational duals remove most core--shell pairs. For a
retained dual, every incidence variable has total coefficient at most one,
while the weighted right side exceeds the global incidence budget. All
coefficients are checked as rational numbers.

The remaining pairs are reduced to integral row-size and side-count
states. The production searches are deterministic and solver-free. For
$m=58,\ldots,64$, the core is $K_{8,8}$ with $64-m$ cross-edges
deleted. Its two sides control the row domains and core-cover conditions.
Table~\ref{tab:o26} gives the proof-facing accounting.

\begin{table}[ht]
\centering
\caption{Order-$26$, degree-nine proof ledger.}
\label{tab:o26}
\scriptsize
\begin{tabular}{@{}c r r L{4.1cm} L{3.2cm}@{}}
\toprule
$m$&cores&shells&primary final exclusion&independent audit\\
\midrule
56,57&431&18&minimum-cover and support classification&
6,804 support pairs reconstructed\\
58&50&55&row states and final shell-triangle lemma&
5,056 Boolean states\\
59&20&85&two-hub finite recurrence&11,943 UNSAT\\
60&10&152&two-hub recurrence and $K_{5,3}+K_1$ lemma&
23,157 UNSAT\\
61&4&393&side counts and set-system recurrence&29,155 UNSAT\\
62&2&1,147&side counts and orbit recurrence&38,342 UNSAT\\
63&1&2,245&cover-13 lemma and orbit recurrence&44,504 UNSAT\\
64&1&7,454&three-hub lemma and orbit recurrence&
101,880 CaDiCaL UNSAT\\
\bottomrule
\end{tabular}
\end{table}

At $m=63$, exact duals leave 334 pairs, scalar conditions leave thirteen,
and the orbit recurrence rejects all 10,715 surviving row states. At
$m=64$, exact duals leave 562 pairs from 7,454. The scalar replay checks
101,880 states and leaves 10,639. A human three-hub argument excludes four
states. In each, the three hubs $c,u,w$ form a path with center $c$.
Writing the sides of the core $K_{8,8}$ as $X,Y$, the scalar assignments
give, up to exchanging the sides,
\[
 |D_c\cap X|=1,\quad |D_c\cap Y|=6,\quad
 D_u,D_w\subset X,\quad |D_u|=|D_w|=7.
\]
If $D_c\cap X=\{x_0\}$ and $S=D_c\cap Y$, the two hub-edge cover
conditions force
\[
 D_u=D_w=X\setminus\{x_0\}.
\]
Every row adjacent to all three hubs then equals
$(Y\setminus S)\cup\{x_0\}$. The six vertices of $S$ still need a
second incidence, but the four states have only $4,0,1,2$ remaining row
positions. This is impossible.

A solver-free $S_8\times S_8$ orbit recurrence excludes the other
10,635 states after 3,883,026 orbit representatives and 299,684,034
demand states. Thus every generated pair is rejected. Independent Boolean
reconstruction is an audit of this last conclusion, not its proof premise.
\end{proof}

\section{The 27-vertex terminal}\label{sec:o27}

\begin{theorem}\label{thm:o27}
The family $\cP_3(27)$ is empty.
\end{theorem}

\begin{proof}
Let $H\in\cF(3,27)$. Lemma~\ref{lem:density} gives $e(H)\le135$.
A vertex of degree at most eight has at least 18 target nonneighbors,
which induce a forbidden member of the two-layer family. Hence
$\delta(H)\ge9$, while averaging gives $\delta(H)\le10$.

Lemma~\ref{lem:two-row} excludes $\delta(H)=9$. If $\delta(H)=10$, then
$H$ is 10-regular and has 135 edges. Proposition~\ref{prop:d10}
excludes this branch.
\end{proof}

\begin{proposition}[Regular endpoint]\label{prop:d10}
There is no 10-regular graph $H\in\cF(3,27)$.
\end{proposition}

\begin{proof}
Choose $v\in V(H)$ and put
\[
 A=N_H(v),\quad B=V(H)\setminus(A\cup\{v\}),\quad
 F=\comp{H[A]},\quad L=\comp{H[B]}.
\]
Thus $|A|=10$, $|B|=16$. For $b\in B$, let
$C_b=N_H(b)\cap A$. Regularity gives
\begin{equation}\label{eq:column-size}
 |C_b|=d_L(b)-5.
\end{equation}
The graph $F$ has independence number at most seven, so
$\tau(F)\ge3$. For every edge $bb'\in E(L)$, the union
$C_b\cup C_{b'}$ covers $F$; otherwise an uncovered edge of $F$ and
$bb'$ form an independent four-set in $H$. Consequently
\begin{equation}\label{eq:degree-sum}
 d_L(b)+d_L(b')\ge13\qquad(bb'\in E(L)).
\end{equation}

Independent canonical generation produces 332 eligible core types, with
edge-level counts
\[
179,80,39,17,9,4,2,1,1
\]
for $e(L)=56,\ldots,64$. The minimum endpoint-degree-sum profile is
\[
\begin{array}{c|rrrrrrr}
\text{minimum sum}&10&11&12&13&14&15&16\\ \hline
\text{cores}&19&102&161&38&10&1&1.
\end{array}
\]
Thus \eqref{eq:degree-sum} rejects 282 cores and leaves 50.

For each survivor, a deterministic CNF encodes a relaxation of the
configuration. It contains exact degrees and edge counts, every split of
an independent four-set, every mixed target $K_9$, the full-color
11-set cap, and the local deletion inequalities. Constraints omitted
from the formula only make it weaker, so an actual configuration would
satisfy the CNF.

The 50 formulas are UNSAT. Their LRAT refutations are independently
replayed. The semantic verifier reconstructs all 332 cores, recomputes the
50 indices, rebuilds each formula without importing the generator, compares
clause multisets, checks every unary totalizer shape, and replays every
LRAT proof. Its aggregate receipt is
\begin{verbatim}
cases_reconstructed=332
cnfs_verified=50
clauses_verified=9100515
totalizer_merges_verified=79574
totalizer_shapes_verified=83
totalizer_states_verified=108839
lrat_proofs_replayed=50
status=PASS
\end{verbatim}
This excludes all 50 survivors and proves the proposition.
\end{proof}

\section{The full-color bridge}\label{sec:bridge}

The terminal statements interact more strongly than scalar averaging
suggests.

\begin{lemma}\label{lem:bridge}
Every $H\in\cF(4,37)$ has at least 192 edges.
\end{lemma}

\begin{proof}
Suppose $e(H)\le191$. If a vertex had degree at most nine, at least 27
of its nonneighbors would induce a member of $\cF(3,27)$, contrary to
Theorem~\ref{thm:o27}. Hence $\delta(H)\ge10$. Since
$2e(H)/37<11$, there is a degree-ten vertex $v$.

Put
\[
 A=N_H(v),\qquad
 B=V(H)\setminus(A\cup\{v\}),\qquad
 M=45-e(H[A]).
\]
Thus $|A|=10$ and $|B|=26$. The 11-set
$\{v\}\cup A$ has $10+45-M$ target edges. By
$D_9(11)=39$,
\begin{equation}\label{eq:M-lower}
 M\ge16.
\end{equation}

For $a\in A$, put $D_a=N_H(a)\cap B$. If
$F=\comp{H[A]}$, minimum degree gives
\[
 10\le d_H(a)=1+9-d_F(a)+|D_a|,
\]
so $|D_a|\ge d_F(a)$. Therefore
\begin{equation}\label{eq:cross-M}
 e_H(A,B)=\sum_{a\in A}|D_a|\ge2M.
\end{equation}
The graph $H[B]$ belongs to $\cF(3,26)$ because $v$ is target-anticomplete
to $B$. Theorem~\ref{thm:o26} gives $e(H[B])\ge121$. Keeping the
missing-edge variable throughout,
\[
\begin{aligned}
 e(H)
 &=10+(45-M)+e_H(A,B)+e(H[B])\\
 &\ge10+(45-M)+2M+121\\
 &=176+M\\
 &\ge192,
\end{aligned}
\]
contrary to the assumption. This proves the lemma.
\end{proof}

\begin{remark}
It is invalid to replace the last display by
$10+29+32+121$: the inequality $e(H[A])\le29$ points in the wrong
direction for a lower bound. The variable $M$ is needed because the
cross-edge gain $2M$ compensates for the shell term $45-M$.
\end{remark}

\section{Completion of the proof}\label{sec:final}

Insert Theorems~\ref{thm:o26} and \ref{thm:o27} and
Lemma~\ref{lem:bridge} into Proposition~\ref{prop:recurrence}. Along the
degree-eight diagonal one obtains
\begin{equation}\label{eq:diagonal}
\begin{array}{c|ccccc}
(a,n)&(4,37)&(5,46)&(6,55)&(7,64)&(8,73)\\ \hline
B(a,n)&192&227&264&301&338.
\end{array}
\end{equation}
Table~\ref{tab:margins} gives all outer packing margins. The symbol
$\emptycert$ denotes a certified empty residual family.

\begin{table}[ht]
\centering
\caption{Margins $B(8-j,81-d-9j)-E_{d,j}$.}
\label{tab:margins}
\footnotesize
\begin{tabular}{c@{\qquad}ccccc}
\toprule
$d$&$j=0$&$j=1$&$j=2$&$j=3$&$j=4$\\
\midrule
0&79&82&$\emptycert$&$\emptycert$&$\emptycert$\\
1&66&69&$\emptycert$&$\emptycert$&$\emptycert$\\
2&56&55&72&$\emptycert$&$\emptycert$\\
3&43&42&45&$\emptycert$&$\emptycert$\\
4&35&34&35&$\emptycert$&$\emptycert$\\
5&24&23&22&$\emptycert$&$\emptycert$\\
6&18&17&16&17&$\emptycert$\\
7&9&8&7&6&$\emptycert$\\
8&5&4&3&2&3\\
\bottomrule
\end{tabular}
\end{table}

\begin{proof}[Proof of Theorem~\ref{thm:main}]
Choose a least color graph $G$ and a minimum-degree vertex $v$. By
\eqref{eq:least}, its degree $d$ lies in $\{0,\ldots,8\}$. Every entry
of the corresponding row in Table~\ref{tab:margins} is positive or
certifies an empty no-next-block family. Applying
\eqref{eq:packing-test} successively forces five disjoint target copies
of $K_9$ in the nonneighbor set of $v$. This contradicts
Lemma~\ref{lem:five-block}. Therefore the assumed coloring does not
exist.
\end{proof}

\section{Certificate semantics and trust boundary}\label{sec:trust}

The proof has four distinct verification layers.
\begin{enumerate}[label=\textup{V\arabic*.},leftmargin=3em]
\item \textbf{Human reduction.} Sections~\ref{sec:colored}--\ref{sec:final}
connect an arbitrary coloring to finite graph and incidence families.
\item \textbf{Catalog completeness.} Canonical graph streams are generated
with nauty's \texttt{geng} \cite{McKayPiperno2014}. Independent programs
check every graph property, count, and graph6 digest used in the paper.
\item \textbf{Finite exclusion.} Rational duals are checked over exact
fractions. Solver-free searches enumerate declared states and use
deterministic recurrences. The order-$27$ branch uses checked LRAT
refutations \cite{CruzFilipeEtAl2017}.
\item \textbf{Independent reconstruction.} Separate source rebuilds the
Boolean semantics, compares generated clauses as multisets, and checks
the catalog classifications. These programs do not treat solver logs as
proofs.
\end{enumerate}

The order-$27$ certificate is a relaxation: it omits some consequences of
the full coloring. This direction is sound because every genuine
configuration satisfies the retained clauses. UNSAT excludes a genuine
configuration; SAT would not by itself construct a coloring.

The order-$26$, level-64 theorem is solver-free. Its independent CaDiCaL
run checks the same 101,880 raw row-size states using a separately written
Boolean encoding and incremental totalizers. It returns 101,880 UNSAT,
zero SAT, and zero UNKNOWN, with classification SHA-256
\texttt{e68dfcaf336ff575e0091a757e60874b81cd79966b418df2e}%
\allowbreak\texttt{6487eb698ade331}.
This audit is deliberately not a premise of Proposition~\ref{prop:o26-finite}.

\section{Reproducibility and data availability}\label{sec:repro}

The release repository contains:
\begin{itemize}[leftmargin=2em]
\item all theorem notes and verifier source;
\item graph and shell generators with pinned catalog hashes;
\item exact rational-dual data and semantic checkers;
\item deterministic solver-free replay programs;
\item four order-$27$ LRAT shards with manifests and checksums;
\item corruption tests and a single dependency-ordered replay entry point;
\item the source and PDF of this paper.
\end{itemize}

The public repository location is
\[
 \text{\url{https://github.com/Robby955/erdos-617-fixed-cases}}.
\]
The fixed-$r=9$ files are under \path{r9/}. The four large certificate
archives are attached to release tag
\texttt{fixed-r9-2026-07-21}. Their names, sizes, and SHA-256 values are
recorded in \path{r9/ASSET_INDEX.json}. The dependency-ordered release
replay passed from the exact public directory layout.

The order-$27$ package currently has four manifest SHA-256 values:
\begingroup
\scriptsize
\begin{verbatim}
f332adb9327c02b2b9964a1d781baae75d27371a83995dc3c1f7393b35c09cba
0df642bb21e52310ec19f479d5ec62fe179a0b075738643bb165baeef7dcda40
7427d8c307411374f332056bdd3638ff26c9d70c6c1283458d87f6313fd1bc2f
e1cc2b4b6f2643d16bc3ba64d50e7e1e9cf73d7043a19f632571bf328c158d1d
\end{verbatim}
\endgroup
The release verifier also rejects modified manifests, same-size formula
mutations, truncated formulas, false source hashes, and truncated LRAT
proofs.

\section{AI-use disclosure}\label{sec:ai}

OpenAI GPT 5.6 Sol was used substantially for proof exploration,
intermediate counterexample searches, adversarial criticism, verification
programs, audit planning, research-note organization, and manuscript
drafting. This use was not limited to copy editing.

Model agreement is not treated as mathematical review. The retained claims
are tied to explicit semantics, exact calculations, deterministic finite
searches or replayed proof certificates, independent reconstructions,
hashes, and corruption tests. The author selected the proof program,
decided which arguments and computations were retained, checked the final
implication chain, and assumes responsibility for every claim and error.

\section{Scope and nonclaims}\label{sec:nonclaims}

Theorem~\ref{thm:main} is a fixed-parameter result. It does not imply the
case $r=10$ and does not settle Erd\H{o}s Problem 617 for arbitrary $r$.
No Lean formalization is claimed. At the date on the title page, the
preprint has not received external mathematical review.

\appendix
\section{Dependency map}\label{app:dependencies}

\begin{table}[ht]
\centering
\caption{Logical and computational dependencies.}
\footnotesize
\begin{tabular}{@{}L{3.25cm} L{4.7cm} L{3.7cm}@{}}
\toprule
result&human dependency&finite dependency\\
\midrule
$P_3(26)\ge121$&degree split; core--shell incidence lemmas&
core levels $56$--$64$, exact duals, solver-free replays\\
$\cP_3(27)=\varnothing$&two-row lemma; regularity; degree-sum 13&
332-core catalog and 50 LRAT refutations\\
\addlinespace[0.3em]
$P_4(37)\ge192$&Theorems~\ref{thm:o26}, \ref{thm:o27};
the $M$-variable argument&integer arithmetic verifier\\
all 45 outer cells&recursive bound; representative selection&
exact integer recurrence replay\\
Theorem~\ref{thm:main}&least-color reduction; five-block exclusion&
the preceding dependency closure\\
\bottomrule
\end{tabular}
\end{table}

Canonical proof notes and replay commands are listed in
\path{output/ERDOS_617_R9_RELEASE_HANDOFF.md}. That handoff records
the exact release commits, artifact locations, tool versions, and
nonclaims. Historical checkpoint notes are not authoritative after the
release commit.

\enlargethispage{10\baselineskip}
\begingroup
\footnotesize
\bibliographystyle{amsplain}
\bibliography{references}
\endgroup

\end{document}
